\section{Assault under providence}
\subsection{Question 114: The tempter's intention and God's government}\label{q114}
Aquinas treats demonic assault by distinguishing the evil intention from the providential government under which the act occurs. In article~1 malice belongs to the demons, who seek to obstruct human salvation. God does not share that intention. He permits evil and orders its consequences to good. Where Scripture describes an evil spirit as sent in punishment, the justice of the judgment and the spirit's hatred remain different intentions. God's governance does not make him the author of the temptation's wickedness. Aquinas sets divine grace first, and angelic protection second, in explaining why human weakness does not face spiritual enemies without help.\footnote{\ST{I q.~114 a.~1}, corpus and ad~1--2; compare Catechism of the Catholic Church, no.~395.}

Article~2 begins with the meaning of tempting as testing. A trial can seek knowledge for a good purpose; the devil tests in order to discover and exploit a vulnerability. When God tests, the resulting disclosure is for creatures, not an addition to divine knowledge. Demons, by contrast, do need to learn. They observe outward conduct and infer inward dispositions, which only God knows directly. Their intelligence must not be confused with omniscience. Temptation can work through imagination and passion, but it cannot force the will's assent.\footnote{\ST{I q.~114 a.~2}, especially ad~2--3.}

Article~3 denies that every sin follows a particular demonic suggestion. Human appetite and a defective free choice can account for a sinful act without such an intervention. Aquinas calls the devil an indirect cause in relation to the first temptation and the wounded condition inherited from that fall; this is different from being the direct instigator of each subsequent offense. His point preserves human responsibility. One cannot transfer one's voluntary act to a spiritual agent merely by noting that temptation exists.\footnote{\ST{I q.~114 a.~3}.}

Article~4 applies the strict distinction about miracles. No demon can act beyond the order of all created nature. Nevertheless a deceptive sign need not be an unreal event. It can be a real effect of created powers used to promote a false conclusion, or an appearance that misleads the senses. Aquinas's natural examples include the ancient theory of generation from putrefaction; that theory belongs to the explanatory natural philosophy of the text, not to the authority of the doctrinal distinction. What matters to the argument is the difference between arranging natural causes and creating or transcending them. He expressly denies that a demon can truly restore a dead person to life by its natural power.\footnote{\ST{I q.~114 a.~4}, corpus and ad~1--2; I q.~110 a.~4. For the causal distinction Aquinas receives, see Augustine, \emph{De Trinitate} III.8.13--9.17.}

Article~5 considers whether a defeated tempter can return. Aquinas reports opinions that defeat bars all further temptation, or bars further temptation of that person. He judges a temporary withdrawal from that person more probable, taking Luke 4:13 seriously: the withdrawal is for a time. The article offers no schedule and no guarantee of permanent immunity after a victory. Its conclusion combines God's merciful limit upon trial with the tempter's strategic retreat. The permission remains God's; neither success nor failure gives a demon independent jurisdiction.\footnote{\ST{I q.~114 a.~5}.}

\angeldossier{I q.~114, aa.~1--5}{Church teaching: Satan's power is finite and permitted under providence, Catechism no.~395. Thomist position: the account of temptation and created powers. Disputed: the duration of withdrawal after defeat.}{Ephesians 6:12; 1~Thessalonians 3:5; Matthew 4:11; Luke 4:13. Aquinas also invokes \emph{De ecclesiasticis dogmatibus}, Augustine, Origen, Ambrose, and the anonymous \emph{Opus imperfectum}.}{Article~5 reports opinions excluding all further temptation, or further temptation of the same person; Aquinas judges temporary withdrawal more probable.}
